\documentclass[11pt, a4paper]{article}
\usepackage{polyglossia}
\setdefaultlanguage{czech}
\usepackage[margin=1in]{geometry}

\makeatletter
\def\input@path{{../configs/}{./}}
\makeatother

\usepackage{math}

\begin{document}
\title{\bb{Semestrální test}}
\author{Jakub Adamec\\B2M01TIK}

\maketitle
Maximální zisk 20 bodů. K zápočtu je třeba získat alespoň 15 bodů. Počítejte algoritmicky, výsledky nehádejte.
\begin{enumerate}
    \item \relax [4 body] Je dána cyklická grupa $\mathbb{Z}_{29}^*$.
    \begin{enumerate}
        \item \relax [2 body] Nalezněte generátor této grupy. Jaká je pravděpodobnost, že při náhodné volbě prvku z grupy $\mathbb{Z}_{29}^*$ trefíte generátor?
        \item \relax [2 body] Najděte všechna řešení rovnice $x^{12} = 1$ v grupě $\mathbb{Z}_{29}^*$.
    \end{enumerate}

    \item \relax [5 bodů] Lineární kód $K$ délky 5 nad $\mathbb{Z}_5$ má generující matici 
    $\mathbf{G} = \begin{pmatrix} 
    2 & 1 & 1 & 3 & 4 \\ 
    3 & 0 & 1 & 3 & 0 \\ 
    4 & 1 & 1 & 4 & 0 
    \end{pmatrix}$.
    \begin{enumerate}
        \item \relax [1 bod] Najděte nějakou kontrolní matici $\mathbf{H}$ kódu $K$.
        \item \relax [2 body] Ověřte, zda je slovo $\bar{w} = (4\ 2\ 2\ 0\ 2)$ bez chyb, případně jednu chybu opravte (pokud to lze).
        \item \relax [2 body] Kolik chyb kód $K$ opravuje?
    \end{enumerate}

    \item \relax [6 bodů] Je dán faktorový okruh $A = \mathbb{Z}_3[x]/q(x)$, kde $q(x) = x^3 + x^2 + 2$.
    \begin{enumerate}
        \item \relax [2 body] Ověřte, že okruh $A$ je těleso. Popište jeho prvky jako polynomy v proměnné $\alpha$, odvoďte potřebná přepisovací pravidla pro násobení dvou prvků.
        \item \relax [2 body] Najděte $(\alpha^2 + 1)^{-1}$ v tělese $A$.
        \item \relax [2 body] Určete řád prvku $\alpha$ v grupě $A^*$.
    \end{enumerate}

    \item \relax [5 bodů] Cyklický kód $K$ délky 6 nad $\mathbb{Z}_5$ má generující polynom $g(z) = z^2 + 4z + 1$.
    \begin{enumerate}
        \item \relax [1 bod] Ověřte, že polynom $g(z) = z^2 + 4z + 1$ může generovat cyklický kód délky 6 nad $\mathbb{Z}_5$. Spočtěte kontrolní polynom $h(z)$ kódu $K$.
        \item \relax [2 body] Pomocí $g(z)$ zakódujte systematicky informaci $\bar{a} = (0\ 1\ 2\ 3)$.
        \item \relax [2 body] Ověřte pomocí kontrolního polynomu $h(z)$, zda je slovo $\bar{w} = (1\ 0\ 2\ 0\ 0\ 1)$ bez chyb, případně jednu chybu opravte (pokud to lze). Podle jakých pravidel násobíte?
    \end{enumerate}
\end{enumerate}

\end{document}
